%%%%%%%%%%%%%%%%%%%%%%%%%%%%%%%%%%%%%%%%%%%%%%%%%%%%%%%%%%%%%%%%%%%%%%%%%%%%%%%%
%  Policy Evaluation Lab -- Lab 5: A case study, an energy price shock
%  Friday 20 November 2026. Same format as Labs 2-4: the case, Table 1 and
%  the questions; groups of 2-3 discuss; review together at the end.
%  Country A and all its numbers are INVENTED (stated on the slides).
%  The review contains only the lecturer's conclusions, kept simple, plus an
%  evaluation slide and a 2022 comparison checked against Eurostat euro
%  indicators 130/2022 (Oct 2022: EA 10.6%, IT 12.6%, energy 41.5%) and the
%  ECB Annual Report 2022 (support ~1.9% of GDP, "largely non-targeted",
%  "about half" affected the marginal cost of energy; negotiated wages +2.8%
%  in 2022; inflation 8.4%; job retention schemes fell; unemployment 6.7% at
%  end-2022); DFR path from the ECB key rates table.
%  Two PDFs are built from this file:
%   - students:  pdflatex lab-05-oil-shock.tex           (case only)
%   - full:      pdflatex "\def\pelReview{}%%%%%%%%%%%%%%%%%%%%%%%%%%%%%%%%%%%%%%%%%%%%%%%%%%%%%%%%%%%%%%%%%%%%%%%%%%%%%%%%
%  Policy Evaluation Lab -- Lab 5: A case study, an energy price shock
%  Friday 20 November 2026. Same format as Labs 2-4: the case, Table 1 and
%  the questions; groups of 2-3 discuss; review together at the end.
%  Country A and all its numbers are INVENTED (stated on the slides).
%  The review contains only the lecturer's conclusions, kept simple, plus an
%  evaluation slide and a 2022 comparison checked against Eurostat euro
%  indicators 130/2022 (Oct 2022: EA 10.6%, IT 12.6%, energy 41.5%) and the
%  ECB Annual Report 2022 (support ~1.9% of GDP, "largely non-targeted",
%  "about half" affected the marginal cost of energy; negotiated wages +2.8%
%  in 2022; inflation 8.4%; job retention schemes fell; unemployment 6.7% at
%  end-2022); DFR path from the ECB key rates table.
%  Two PDFs are built from this file:
%   - students:  pdflatex lab-05-oil-shock.tex           (case only)
%   - full:      pdflatex "\def\pelReview{}%%%%%%%%%%%%%%%%%%%%%%%%%%%%%%%%%%%%%%%%%%%%%%%%%%%%%%%%%%%%%%%%%%%%%%%%%%%%%%%%
%  Policy Evaluation Lab -- Lab 5: A case study, an energy price shock
%  Friday 20 November 2026. Same format as Labs 2-4: the case, Table 1 and
%  the questions; groups of 2-3 discuss; review together at the end.
%  Country A and all its numbers are INVENTED (stated on the slides).
%  The review contains only the lecturer's conclusions, kept simple, plus an
%  evaluation slide and a 2022 comparison checked against Eurostat euro
%  indicators 130/2022 (Oct 2022: EA 10.6%, IT 12.6%, energy 41.5%) and the
%  ECB Annual Report 2022 (support ~1.9% of GDP, "largely non-targeted",
%  "about half" affected the marginal cost of energy; negotiated wages +2.8%
%  in 2022; inflation 8.4%; job retention schemes fell; unemployment 6.7% at
%  end-2022); DFR path from the ECB key rates table.
%  Two PDFs are built from this file:
%   - students:  pdflatex lab-05-oil-shock.tex           (case only)
%   - full:      pdflatex "\def\pelReview{}%%%%%%%%%%%%%%%%%%%%%%%%%%%%%%%%%%%%%%%%%%%%%%%%%%%%%%%%%%%%%%%%%%%%%%%%%%%%%%%%
%  Policy Evaluation Lab -- Lab 5: A case study, an energy price shock
%  Friday 20 November 2026. Same format as Labs 2-4: the case, Table 1 and
%  the questions; groups of 2-3 discuss; review together at the end.
%  Country A and all its numbers are INVENTED (stated on the slides).
%  The review contains only the lecturer's conclusions, kept simple, plus an
%  evaluation slide and a 2022 comparison checked against Eurostat euro
%  indicators 130/2022 (Oct 2022: EA 10.6%, IT 12.6%, energy 41.5%) and the
%  ECB Annual Report 2022 (support ~1.9% of GDP, "largely non-targeted",
%  "about half" affected the marginal cost of energy; negotiated wages +2.8%
%  in 2022; inflation 8.4%; job retention schemes fell; unemployment 6.7% at
%  end-2022); DFR path from the ECB key rates table.
%  Two PDFs are built from this file:
%   - students:  pdflatex lab-05-oil-shock.tex           (case only)
%   - full:      pdflatex "\def\pelReview{}\input{lab-05-oil-shock.tex}"
%%%%%%%%%%%%%%%%%%%%%%%%%%%%%%%%%%%%%%%%%%%%%%%%%%%%%%%%%%%%%%%%%%%%%%%%%%%%%%%%

\documentclass[aspectratio=169,xcolor=table]{beamer}
\input{preamble-pel}
\newcolumntype{L}[1]{>{\raggedright\arraybackslash}p{#1}}
\tikzset{
  ax/.style={-{Stealth[length=2mm]}, thick, pelText},
  lbl/.style={font=\scriptsize},
  guide/.style={densely dotted, pelMuted},
}

\title[Policy Evaluation Lab]{Policy Evaluation Lab}
\subtitle[Lab 5]{Lab 5: A case study, an energy price shock}
\date[\pelYear]{November, 2026}

\begin{document}

\frame{\titlepage}

\begin{frame}{So, today we are going to evaluate a case study}
\be
  \item We look at a case study: a country hit by an energy shock
  \item Form groups of \textbf{2--3 people} at most
  \item Study what the questions are asking, and look at the data
  \item Formulate some opinions, and the policies you would recommend
  \item Present your results; at the end we review the case together
\ee
\end{frame}

\begin{frame}{Today's lab: a case study}
\footnotesize
\textbf{Country A} is a medium-sized, high-income economy with its own currency and an independent central bank, which targets inflation of \textbf{2\%}. The country has no oil or gas of its own: almost all its energy is imported. Its manufacturing sector is large and energy-intensive (steel, chemicals, ceramics, transport). Last year the economy grew by 1.8\%, unemployment was at 7\%, which the central bank considers its natural rate, and inflation was on target. Public debt is high, at \textbf{95\% of GDP}, with a deficit of 2.5\%.

\vspace{0.15cm}
In January, a war between two oil-producing countries disrupts world supply. The price of oil \textbf{rises by 60\%} in two months, and markets expect it to \textbf{stay high for years}. Energy bills for households and firms jump. The main unions ask for wage increases to recover the lost purchasing power, and several energy-intensive firms announce temporary plant closures. The government and the central bank must decide how to respond. The first estimates for this year are reported in Table 1.
\pelTableNote{Country A and all the numbers in this case are invented for the exercise.}
\end{frame}

\begin{frame}{Table 1: Country A before and after the shock}
\centering\footnotesize
\begin{tabular}{lrr}
\toprule
 & \textbf{Last year} & \textbf{This year (first estimate)}\\
\midrule
Real GDP growth & 1.8\% & 0.6\% \\
Unemployment rate (natural rate: 7\%) & 7.0\% & 7.2\% \\
Inflation (all items) & 2.0\% & 4.8\% \\
Inflation excluding energy & 2.0\% & 2.6\% \\
Energy prices for consumers & $+1\%$ & $+25\%$ \\
Central bank policy rate & 3.0\% & 3.0\% \\
Government deficit (\% of GDP) & 2.5 & 2.5 \\
Government debt (\% of GDP) & 95 & 95 \\
Energy imports (\% of GDP) & 4 & 6 \\
Consumer confidence (long-run average = 100) & 101 & 88 \\
\bottomrule
\end{tabular}
\pelTableNote{All numbers invented for the exercise. Energy is about 10\% of the consumer basket. ``This year'': before any policy response.}
\end{frame}

\begin{frame}{The questions}
\small
\begin{block}{}
Based on the information provided, analyse the nature of the shock and its expected effects on output, inflation and unemployment, in the short and in the medium run.
\end{block}
\vspace{0.05cm}
\begin{block}{}
Illustrate which policies the government and the central bank could adopt in response, discussing their advantages, their costs and the constraints they face.
\end{block}
\vspace{0.05cm}
\begin{block}{}
In addition, illustrate how you would evaluate, after two years, whether the measures adopted were effective, and which data would be needed.
\end{block}
\end{frame}

\ifdefined\pelReview

\section{Let's review it together}

\begin{frame}{What we expect}
\bi
  \item \textbf{Higher prices}: Country A imports its energy, so the shock comes straight into prices
  \item \textbf{Firms with higher costs}, above all in the energy-intensive sectors: possible downsizing
  \item \textbf{Poorer households}: real wages fall
  \item The shock may be \textbf{permanent}: the country is poorer, and the economy suffers anyway
\ei
\vspace{0.2cm}
\begin{block}{}
\centering The question is how to \textbf{absorb} the shock and \textbf{redistribute} its effects, safeguarding the poorest households
\end{block}
\end{frame}

\begin{frame}{The government}
\bi
  \item \textbf{Energy subsidies} to firms and households: yes, but \textbf{limited}, \textbf{temporary} and \textbf{targeted} at the poorest households
  \item \textbf{Furlough schemes} for firms hit by the costs: they protect jobs while it is unclear how long the shock will last, with an \textbf{end date}
  \item \textbf{Wages}: raising them to recover all the lost purchasing power can start a \textbf{wage-price spiral}
\ei
\end{frame}

\begin{frame}{The central bank}
\bi
  \item \textbf{Raise interest rates} to slow down inflation, above all if wages and other prices start to follow energy prices (the spiral)
  \item \textbf{Coordinated} with a specific and accurate fiscal policy: support that is targeted, so that it does not push demand and inflation up while rates are rising
\ei
\end{frame}

\begin{frame}{What we conclude}
\begin{block}{}
A permanent energy shock makes the country poorer: no policy can undo it. Policy decides \textbf{who bears the loss} and \textbf{how the economy adjusts}:
\bi
  \item temporary, targeted support for the poorest households
  \item furlough schemes with an end date
  \item no full wage indexation, to avoid a spiral
  \item higher interest rates, coordinated with fiscal policy
\ei
\end{block}
\end{frame}

\begin{frame}{How would we evaluate it?}
\bi
  \item \textbf{The effect of the shock}: in the following years, compare Country A with similar countries that were \textbf{less exposed} (smaller energy importers), before and after the shock
  \item But Country A differs in two ways: it was hit harder \textbf{and} it adopted its own policies. That comparison mixes the two
  \item \textbf{The effect of the policies}: compare who received them with who did not
  \bi
    \item households just above and just below the threshold for the subsidies
    \item firms that used the furlough scheme versus similar firms that did not
    \item or countries hit just as hard that chose \textbf{different} policies
  \ei
\ei
\end{frame}

\begin{frame}{Real life: the 2022 energy shock in the euro area}
\footnotesize
\renewcommand{\arraystretch}{1.05}
\begin{tabular}{L{4.0cm}L{8.6cm}}
\toprule
\textbf{What we said} & \textbf{What happened} \\
\midrule
Higher prices & Inflation reached \textbf{10.6\%} in October 2022 (Italy 12.6\%); energy prices $+41.5\%$ in a year \\
Limited, targeted subsidies & About \textbf{1.9\% of GDP} in 2022, \textbf{largely non-targeted}; about half lowered the price of energy itself \\
Furlough schemes & Not the main tool: their use fell in 2022; unemployment at a historic low (6.7\%) \\
No full wage indexation & Negotiated wages $+2.8\%$ in 2022 against inflation of 8.4\%: real wages fell, no spiral at first \\
Raise interest rates & The ECB raised rates from July 2022: deposit rate from $-0.50\%$ to \textbf{4.00\%} by September 2023 \\
Coordination & Not quite: governments expanded support while the ECB was tightening \\
\bottomrule
\end{tabular}
\pelTableNote{Sources: Eurostat, euro indicators 130/2022 (17 November 2022); ECB Annual Report 2022; ECB, Key ECB interest rates.}
\end{frame}

\fi

\pelQuestions

\end{document}
"
%%%%%%%%%%%%%%%%%%%%%%%%%%%%%%%%%%%%%%%%%%%%%%%%%%%%%%%%%%%%%%%%%%%%%%%%%%%%%%%%

\documentclass[aspectratio=169,xcolor=table]{beamer}
\input{preamble-pel}
\newcolumntype{L}[1]{>{\raggedright\arraybackslash}p{#1}}
\tikzset{
  ax/.style={-{Stealth[length=2mm]}, thick, pelText},
  lbl/.style={font=\scriptsize},
  guide/.style={densely dotted, pelMuted},
}

\title[Policy Evaluation Lab]{Policy Evaluation Lab}
\subtitle[Lab 5]{Lab 5: A case study, an energy price shock}
\date[\pelYear]{November, 2026}

\begin{document}

\frame{\titlepage}

\begin{frame}{So, today we are going to evaluate a case study}
\be
  \item We look at a case study: a country hit by an energy shock
  \item Form groups of \textbf{2--3 people} at most
  \item Study what the questions are asking, and look at the data
  \item Formulate some opinions, and the policies you would recommend
  \item Present your results; at the end we review the case together
\ee
\end{frame}

\begin{frame}{Today's lab: a case study}
\footnotesize
\textbf{Country A} is a medium-sized, high-income economy with its own currency and an independent central bank, which targets inflation of \textbf{2\%}. The country has no oil or gas of its own: almost all its energy is imported. Its manufacturing sector is large and energy-intensive (steel, chemicals, ceramics, transport). Last year the economy grew by 1.8\%, unemployment was at 7\%, which the central bank considers its natural rate, and inflation was on target. Public debt is high, at \textbf{95\% of GDP}, with a deficit of 2.5\%.

\vspace{0.15cm}
In January, a war between two oil-producing countries disrupts world supply. The price of oil \textbf{rises by 60\%} in two months, and markets expect it to \textbf{stay high for years}. Energy bills for households and firms jump. The main unions ask for wage increases to recover the lost purchasing power, and several energy-intensive firms announce temporary plant closures. The government and the central bank must decide how to respond. The first estimates for this year are reported in Table 1.
\pelTableNote{Country A and all the numbers in this case are invented for the exercise.}
\end{frame}

\begin{frame}{Table 1: Country A before and after the shock}
\centering\footnotesize
\begin{tabular}{lrr}
\toprule
 & \textbf{Last year} & \textbf{This year (first estimate)}\\
\midrule
Real GDP growth & 1.8\% & 0.6\% \\
Unemployment rate (natural rate: 7\%) & 7.0\% & 7.2\% \\
Inflation (all items) & 2.0\% & 4.8\% \\
Inflation excluding energy & 2.0\% & 2.6\% \\
Energy prices for consumers & $+1\%$ & $+25\%$ \\
Central bank policy rate & 3.0\% & 3.0\% \\
Government deficit (\% of GDP) & 2.5 & 2.5 \\
Government debt (\% of GDP) & 95 & 95 \\
Energy imports (\% of GDP) & 4 & 6 \\
Consumer confidence (long-run average = 100) & 101 & 88 \\
\bottomrule
\end{tabular}
\pelTableNote{All numbers invented for the exercise. Energy is about 10\% of the consumer basket. ``This year'': before any policy response.}
\end{frame}

\begin{frame}{The questions}
\small
\begin{block}{}
Based on the information provided, analyse the nature of the shock and its expected effects on output, inflation and unemployment, in the short and in the medium run.
\end{block}
\vspace{0.05cm}
\begin{block}{}
Illustrate which policies the government and the central bank could adopt in response, discussing their advantages, their costs and the constraints they face.
\end{block}
\vspace{0.05cm}
\begin{block}{}
In addition, illustrate how you would evaluate, after two years, whether the measures adopted were effective, and which data would be needed.
\end{block}
\end{frame}

\ifdefined\pelReview

\section{Let's review it together}

\begin{frame}{What we expect}
\bi
  \item \textbf{Higher prices}: Country A imports its energy, so the shock comes straight into prices
  \item \textbf{Firms with higher costs}, above all in the energy-intensive sectors: possible downsizing
  \item \textbf{Poorer households}: real wages fall
  \item The shock may be \textbf{permanent}: the country is poorer, and the economy suffers anyway
\ei
\vspace{0.2cm}
\begin{block}{}
\centering The question is how to \textbf{absorb} the shock and \textbf{redistribute} its effects, safeguarding the poorest households
\end{block}
\end{frame}

\begin{frame}{The government}
\bi
  \item \textbf{Energy subsidies} to firms and households: yes, but \textbf{limited}, \textbf{temporary} and \textbf{targeted} at the poorest households
  \item \textbf{Furlough schemes} for firms hit by the costs: they protect jobs while it is unclear how long the shock will last, with an \textbf{end date}
  \item \textbf{Wages}: raising them to recover all the lost purchasing power can start a \textbf{wage-price spiral}
\ei
\end{frame}

\begin{frame}{The central bank}
\bi
  \item \textbf{Raise interest rates} to slow down inflation, above all if wages and other prices start to follow energy prices (the spiral)
  \item \textbf{Coordinated} with a specific and accurate fiscal policy: support that is targeted, so that it does not push demand and inflation up while rates are rising
\ei
\end{frame}

\begin{frame}{What we conclude}
\begin{block}{}
A permanent energy shock makes the country poorer: no policy can undo it. Policy decides \textbf{who bears the loss} and \textbf{how the economy adjusts}:
\bi
  \item temporary, targeted support for the poorest households
  \item furlough schemes with an end date
  \item no full wage indexation, to avoid a spiral
  \item higher interest rates, coordinated with fiscal policy
\ei
\end{block}
\end{frame}

\begin{frame}{How would we evaluate it?}
\bi
  \item \textbf{The effect of the shock}: in the following years, compare Country A with similar countries that were \textbf{less exposed} (smaller energy importers), before and after the shock
  \item But Country A differs in two ways: it was hit harder \textbf{and} it adopted its own policies. That comparison mixes the two
  \item \textbf{The effect of the policies}: compare who received them with who did not
  \bi
    \item households just above and just below the threshold for the subsidies
    \item firms that used the furlough scheme versus similar firms that did not
    \item or countries hit just as hard that chose \textbf{different} policies
  \ei
\ei
\end{frame}

\begin{frame}{Real life: the 2022 energy shock in the euro area}
\footnotesize
\renewcommand{\arraystretch}{1.05}
\begin{tabular}{L{4.0cm}L{8.6cm}}
\toprule
\textbf{What we said} & \textbf{What happened} \\
\midrule
Higher prices & Inflation reached \textbf{10.6\%} in October 2022 (Italy 12.6\%); energy prices $+41.5\%$ in a year \\
Limited, targeted subsidies & About \textbf{1.9\% of GDP} in 2022, \textbf{largely non-targeted}; about half lowered the price of energy itself \\
Furlough schemes & Not the main tool: their use fell in 2022; unemployment at a historic low (6.7\%) \\
No full wage indexation & Negotiated wages $+2.8\%$ in 2022 against inflation of 8.4\%: real wages fell, no spiral at first \\
Raise interest rates & The ECB raised rates from July 2022: deposit rate from $-0.50\%$ to \textbf{4.00\%} by September 2023 \\
Coordination & Not quite: governments expanded support while the ECB was tightening \\
\bottomrule
\end{tabular}
\pelTableNote{Sources: Eurostat, euro indicators 130/2022 (17 November 2022); ECB Annual Report 2022; ECB, Key ECB interest rates.}
\end{frame}

\fi

\pelQuestions

\end{document}
"
%%%%%%%%%%%%%%%%%%%%%%%%%%%%%%%%%%%%%%%%%%%%%%%%%%%%%%%%%%%%%%%%%%%%%%%%%%%%%%%%

\documentclass[aspectratio=169,xcolor=table]{beamer}
\input{preamble-pel}
\newcolumntype{L}[1]{>{\raggedright\arraybackslash}p{#1}}
\tikzset{
  ax/.style={-{Stealth[length=2mm]}, thick, pelText},
  lbl/.style={font=\scriptsize},
  guide/.style={densely dotted, pelMuted},
}

\title[Policy Evaluation Lab]{Policy Evaluation Lab}
\subtitle[Lab 5]{Lab 5: A case study, an energy price shock}
\date[\pelYear]{November, 2026}

\begin{document}

\frame{\titlepage}

\begin{frame}{So, today we are going to evaluate a case study}
\be
  \item We look at a case study: a country hit by an energy shock
  \item Form groups of \textbf{2--3 people} at most
  \item Study what the questions are asking, and look at the data
  \item Formulate some opinions, and the policies you would recommend
  \item Present your results; at the end we review the case together
\ee
\end{frame}

\begin{frame}{Today's lab: a case study}
\footnotesize
\textbf{Country A} is a medium-sized, high-income economy with its own currency and an independent central bank, which targets inflation of \textbf{2\%}. The country has no oil or gas of its own: almost all its energy is imported. Its manufacturing sector is large and energy-intensive (steel, chemicals, ceramics, transport). Last year the economy grew by 1.8\%, unemployment was at 7\%, which the central bank considers its natural rate, and inflation was on target. Public debt is high, at \textbf{95\% of GDP}, with a deficit of 2.5\%.

\vspace{0.15cm}
In January, a war between two oil-producing countries disrupts world supply. The price of oil \textbf{rises by 60\%} in two months, and markets expect it to \textbf{stay high for years}. Energy bills for households and firms jump. The main unions ask for wage increases to recover the lost purchasing power, and several energy-intensive firms announce temporary plant closures. The government and the central bank must decide how to respond. The first estimates for this year are reported in Table 1.
\pelTableNote{Country A and all the numbers in this case are invented for the exercise.}
\end{frame}

\begin{frame}{Table 1: Country A before and after the shock}
\centering\footnotesize
\begin{tabular}{lrr}
\toprule
 & \textbf{Last year} & \textbf{This year (first estimate)}\\
\midrule
Real GDP growth & 1.8\% & 0.6\% \\
Unemployment rate (natural rate: 7\%) & 7.0\% & 7.2\% \\
Inflation (all items) & 2.0\% & 4.8\% \\
Inflation excluding energy & 2.0\% & 2.6\% \\
Energy prices for consumers & $+1\%$ & $+25\%$ \\
Central bank policy rate & 3.0\% & 3.0\% \\
Government deficit (\% of GDP) & 2.5 & 2.5 \\
Government debt (\% of GDP) & 95 & 95 \\
Energy imports (\% of GDP) & 4 & 6 \\
Consumer confidence (long-run average = 100) & 101 & 88 \\
\bottomrule
\end{tabular}
\pelTableNote{All numbers invented for the exercise. Energy is about 10\% of the consumer basket. ``This year'': before any policy response.}
\end{frame}

\begin{frame}{The questions}
\small
\begin{block}{}
Based on the information provided, analyse the nature of the shock and its expected effects on output, inflation and unemployment, in the short and in the medium run.
\end{block}
\vspace{0.05cm}
\begin{block}{}
Illustrate which policies the government and the central bank could adopt in response, discussing their advantages, their costs and the constraints they face.
\end{block}
\vspace{0.05cm}
\begin{block}{}
In addition, illustrate how you would evaluate, after two years, whether the measures adopted were effective, and which data would be needed.
\end{block}
\end{frame}

\ifdefined\pelReview

\section{Let's review it together}

\begin{frame}{What we expect}
\bi
  \item \textbf{Higher prices}: Country A imports its energy, so the shock comes straight into prices
  \item \textbf{Firms with higher costs}, above all in the energy-intensive sectors: possible downsizing
  \item \textbf{Poorer households}: real wages fall
  \item The shock may be \textbf{permanent}: the country is poorer, and the economy suffers anyway
\ei
\vspace{0.2cm}
\begin{block}{}
\centering The question is how to \textbf{absorb} the shock and \textbf{redistribute} its effects, safeguarding the poorest households
\end{block}
\end{frame}

\begin{frame}{The government}
\bi
  \item \textbf{Energy subsidies} to firms and households: yes, but \textbf{limited}, \textbf{temporary} and \textbf{targeted} at the poorest households
  \item \textbf{Furlough schemes} for firms hit by the costs: they protect jobs while it is unclear how long the shock will last, with an \textbf{end date}
  \item \textbf{Wages}: raising them to recover all the lost purchasing power can start a \textbf{wage-price spiral}
\ei
\end{frame}

\begin{frame}{The central bank}
\bi
  \item \textbf{Raise interest rates} to slow down inflation, above all if wages and other prices start to follow energy prices (the spiral)
  \item \textbf{Coordinated} with a specific and accurate fiscal policy: support that is targeted, so that it does not push demand and inflation up while rates are rising
\ei
\end{frame}

\begin{frame}{What we conclude}
\begin{block}{}
A permanent energy shock makes the country poorer: no policy can undo it. Policy decides \textbf{who bears the loss} and \textbf{how the economy adjusts}:
\bi
  \item temporary, targeted support for the poorest households
  \item furlough schemes with an end date
  \item no full wage indexation, to avoid a spiral
  \item higher interest rates, coordinated with fiscal policy
\ei
\end{block}
\end{frame}

\begin{frame}{How would we evaluate it?}
\bi
  \item \textbf{The effect of the shock}: in the following years, compare Country A with similar countries that were \textbf{less exposed} (smaller energy importers), before and after the shock
  \item But Country A differs in two ways: it was hit harder \textbf{and} it adopted its own policies. That comparison mixes the two
  \item \textbf{The effect of the policies}: compare who received them with who did not
  \bi
    \item households just above and just below the threshold for the subsidies
    \item firms that used the furlough scheme versus similar firms that did not
    \item or countries hit just as hard that chose \textbf{different} policies
  \ei
\ei
\end{frame}

\begin{frame}{Real life: the 2022 energy shock in the euro area}
\footnotesize
\renewcommand{\arraystretch}{1.05}
\begin{tabular}{L{4.0cm}L{8.6cm}}
\toprule
\textbf{What we said} & \textbf{What happened} \\
\midrule
Higher prices & Inflation reached \textbf{10.6\%} in October 2022 (Italy 12.6\%); energy prices $+41.5\%$ in a year \\
Limited, targeted subsidies & About \textbf{1.9\% of GDP} in 2022, \textbf{largely non-targeted}; about half lowered the price of energy itself \\
Furlough schemes & Not the main tool: their use fell in 2022; unemployment at a historic low (6.7\%) \\
No full wage indexation & Negotiated wages $+2.8\%$ in 2022 against inflation of 8.4\%: real wages fell, no spiral at first \\
Raise interest rates & The ECB raised rates from July 2022: deposit rate from $-0.50\%$ to \textbf{4.00\%} by September 2023 \\
Coordination & Not quite: governments expanded support while the ECB was tightening \\
\bottomrule
\end{tabular}
\pelTableNote{Sources: Eurostat, euro indicators 130/2022 (17 November 2022); ECB Annual Report 2022; ECB, Key ECB interest rates.}
\end{frame}

\fi

\pelQuestions

\end{document}
"
%%%%%%%%%%%%%%%%%%%%%%%%%%%%%%%%%%%%%%%%%%%%%%%%%%%%%%%%%%%%%%%%%%%%%%%%%%%%%%%%

\documentclass[aspectratio=169,xcolor=table]{beamer}
\input{preamble-pel}
\newcolumntype{L}[1]{>{\raggedright\arraybackslash}p{#1}}
\tikzset{
  ax/.style={-{Stealth[length=2mm]}, thick, pelText},
  lbl/.style={font=\scriptsize},
  guide/.style={densely dotted, pelMuted},
}

\title[Policy Evaluation Lab]{Policy Evaluation Lab}
\subtitle[Lab 5]{Lab 5: A case study, an energy price shock}
\date[\pelYear]{November, 2026}

\begin{document}

\frame{\titlepage}

\begin{frame}{So, today we are going to evaluate a case study}
\be
  \item We look at a case study: a country hit by an energy shock
  \item Form groups of \textbf{2--3 people} at most
  \item Study what the questions are asking, and look at the data
  \item Formulate some opinions, and the policies you would recommend
  \item Present your results; at the end we review the case together
\ee
\end{frame}

\begin{frame}{Today's lab: a case study}
\footnotesize
\textbf{Country A} is a medium-sized, high-income economy with its own currency and an independent central bank, which targets inflation of \textbf{2\%}. The country has no oil or gas of its own: almost all its energy is imported. Its manufacturing sector is large and energy-intensive (steel, chemicals, ceramics, transport). Last year the economy grew by 1.8\%, unemployment was at 7\%, which the central bank considers its natural rate, and inflation was on target. Public debt is high, at \textbf{95\% of GDP}, with a deficit of 2.5\%.

\vspace{0.15cm}
In January, a war between two oil-producing countries disrupts world supply. The price of oil \textbf{rises by 60\%} in two months, and markets expect it to \textbf{stay high for years}. Energy bills for households and firms jump. The main unions ask for wage increases to recover the lost purchasing power, and several energy-intensive firms announce temporary plant closures. The government and the central bank must decide how to respond. The first estimates for this year are reported in Table 1.
\pelTableNote{Country A and all the numbers in this case are invented for the exercise.}
\end{frame}

\begin{frame}{Table 1: Country A before and after the shock}
\centering\footnotesize
\begin{tabular}{lrr}
\toprule
 & \textbf{Last year} & \textbf{This year (first estimate)}\\
\midrule
Real GDP growth & 1.8\% & 0.6\% \\
Unemployment rate (natural rate: 7\%) & 7.0\% & 7.2\% \\
Inflation (all items) & 2.0\% & 4.8\% \\
Inflation excluding energy & 2.0\% & 2.6\% \\
Energy prices for consumers & $+1\%$ & $+25\%$ \\
Central bank policy rate & 3.0\% & 3.0\% \\
Government deficit (\% of GDP) & 2.5 & 2.5 \\
Government debt (\% of GDP) & 95 & 95 \\
Energy imports (\% of GDP) & 4 & 6 \\
Consumer confidence (long-run average = 100) & 101 & 88 \\
\bottomrule
\end{tabular}
\pelTableNote{All numbers invented for the exercise. Energy is about 10\% of the consumer basket. ``This year'': before any policy response.}
\end{frame}

\begin{frame}{The questions}
\small
\begin{block}{}
Based on the information provided, analyse the nature of the shock and its expected effects on output, inflation and unemployment, in the short and in the medium run.
\end{block}
\vspace{0.05cm}
\begin{block}{}
Illustrate which policies the government and the central bank could adopt in response, discussing their advantages, their costs and the constraints they face.
\end{block}
\vspace{0.05cm}
\begin{block}{}
In addition, illustrate how you would evaluate, after two years, whether the measures adopted were effective, and which data would be needed.
\end{block}
\end{frame}

\ifdefined\pelReview

\section{Let's review it together}

\begin{frame}{What we expect}
\bi
  \item \textbf{Higher prices}: Country A imports its energy, so the shock comes straight into prices
  \item \textbf{Firms with higher costs}, above all in the energy-intensive sectors: possible downsizing
  \item \textbf{Poorer households}: real wages fall
  \item The shock may be \textbf{permanent}: the country is poorer, and the economy suffers anyway
\ei
\vspace{0.2cm}
\begin{block}{}
\centering The question is how to \textbf{absorb} the shock and \textbf{redistribute} its effects, safeguarding the poorest households
\end{block}
\end{frame}

\begin{frame}{The government}
\bi
  \item \textbf{Energy subsidies} to firms and households: yes, but \textbf{limited}, \textbf{temporary} and \textbf{targeted} at the poorest households
  \item \textbf{Furlough schemes} for firms hit by the costs: they protect jobs while it is unclear how long the shock will last, with an \textbf{end date}
  \item \textbf{Wages}: raising them to recover all the lost purchasing power can start a \textbf{wage-price spiral}
\ei
\end{frame}

\begin{frame}{The central bank}
\bi
  \item \textbf{Raise interest rates} to slow down inflation, above all if wages and other prices start to follow energy prices (the spiral)
  \item \textbf{Coordinated} with a specific and accurate fiscal policy: support that is targeted, so that it does not push demand and inflation up while rates are rising
\ei
\end{frame}

\begin{frame}{What we conclude}
\begin{block}{}
A permanent energy shock makes the country poorer: no policy can undo it. Policy decides \textbf{who bears the loss} and \textbf{how the economy adjusts}:
\bi
  \item temporary, targeted support for the poorest households
  \item furlough schemes with an end date
  \item no full wage indexation, to avoid a spiral
  \item higher interest rates, coordinated with fiscal policy
\ei
\end{block}
\end{frame}

\begin{frame}{How would we evaluate it?}
\bi
  \item \textbf{The effect of the shock}: in the following years, compare Country A with similar countries that were \textbf{less exposed} (smaller energy importers), before and after the shock
  \item But Country A differs in two ways: it was hit harder \textbf{and} it adopted its own policies. That comparison mixes the two
  \item \textbf{The effect of the policies}: compare who received them with who did not
  \bi
    \item households just above and just below the threshold for the subsidies
    \item firms that used the furlough scheme versus similar firms that did not
    \item or countries hit just as hard that chose \textbf{different} policies
  \ei
\ei
\end{frame}

\begin{frame}{Real life: the 2022 energy shock in the euro area}
\footnotesize
\renewcommand{\arraystretch}{1.05}
\begin{tabular}{L{4.0cm}L{8.6cm}}
\toprule
\textbf{What we said} & \textbf{What happened} \\
\midrule
Higher prices & Inflation reached \textbf{10.6\%} in October 2022 (Italy 12.6\%); energy prices $+41.5\%$ in a year \\
Limited, targeted subsidies & About \textbf{1.9\% of GDP} in 2022, \textbf{largely non-targeted}; about half lowered the price of energy itself \\
Furlough schemes & Not the main tool: their use fell in 2022; unemployment at a historic low (6.7\%) \\
No full wage indexation & Negotiated wages $+2.8\%$ in 2022 against inflation of 8.4\%: real wages fell, no spiral at first \\
Raise interest rates & The ECB raised rates from July 2022: deposit rate from $-0.50\%$ to \textbf{4.00\%} by September 2023 \\
Coordination & Not quite: governments expanded support while the ECB was tightening \\
\bottomrule
\end{tabular}
\pelTableNote{Sources: Eurostat, euro indicators 130/2022 (17 November 2022); ECB Annual Report 2022; ECB, Key ECB interest rates.}
\end{frame}

\fi

\pelQuestions

\end{document}
